\section{Prerequisite Knowledge and Imported Tools}

\subsection*{Wasserstein gradient flows}
The paper repeatedly uses the idea that the limiting PDE is a gradient flow on the metric space $(\Pcal(\R^D), \Wtwo)$. The relevant imported viewpoint is the Jordan--Kinderlehrer--Otto / Ambrosio--Gigli--Savar\'e framework \cite{jordan1998variational,ambrosio2008gradient}. This is not re-proved in the paper, but it strongly shapes how the PDE is interpreted.

\subsection*{Propagation of chaos}
The step from finite-width SGD to a measure-valued nonlinear limit is an instance of propagation of chaos \cite{sznitman1991topics}. The paper's proof is largely self-contained, but this is the conceptual background.

\subsection*{Measure-valued PDEs}
The distributional dynamics is interpreted weakly in the space of probability measures. A reader should already be comfortable with continuity equations, weak solutions, and measure-valued evolutions.

\subsection*{Classical optimization and SGD scaling}
The paper assumes familiarity with SGD in the small-step-size regime and with the idea that a discrete optimization algorithm can converge to an ODE/PDE under scaling.

\subsection*{Example-specific prerequisites}
The isotropic and anisotropic Gaussian examples use symmetry reduction. The reader should be comfortable with:
\begin{itemize}[leftmargin=1.5em]
\item radial distributions and rotational invariance;
\item reduced coordinates for anisotropic covariance structure;
\item basic sub-Gaussian concentration estimates.
\end{itemize}
